\chapter{The Elm loop and algebraic effects}
\label{ch:elm}

\maya{} organises a whole application around a single idea borrowed from
the Elm language: every program is a pure function from the current
state and an incoming event to the next state. Anything that has to
\emph{happen} in the outside world --- a timer firing, a key being
read, a clipboard being written, an HTTP request --- is returned as a
\emph{description of work to do}, not performed inside the function.
This chapter explains the idea, why it is worth the discipline, and
how \maya{} encodes it.

\begin{aside}
If you have written GUIs before, you have written event handlers ---
short functions hung off buttons and inputs that mutate fields, fire
requests, and redraw. They work. They also tend to accumulate, over
months, into a tangle nobody can hold in their head. Elm's response
was counter-intuitive: \emph{remove} the ability to do things in
handlers, force handlers to return plain data, and have a separate
layer perform the actual work. The architecture sounds restrictive
and turns out to be liberating: there is suddenly one place where
state changes (your \code{update}) and one place where the world is
touched (the runtime). This chapter is the long version of that
sentence.
\end{aside}

\section{What does ``pure'' actually mean?}
\label{sec:elm:purity}

The word ``pure'' gets thrown around in programming circles with the
vigour of a moral verdict; it has a precise meaning that has nothing
to do with virtue.

A function is \textbf{pure} when both of these hold:

\begin{enumerate}[leftmargin=*]
  \item \textbf{Same inputs $\Rightarrow$ same output.} Call
    \code{f(3)} twice and you get the same value both times. The
    function does not consult the clock, the filesystem, or a random
    number generator. It does not look at any global variable.
  \item \textbf{No side effects.} Calling the function does nothing
    observable to the rest of the program: no field is mutated, no
    file written, no byte sent on the wire. The only thing the
    caller can observe is the return value.
\end{enumerate}

Both halves are needed. \code{std::time(nullptr)} satisfies (2) ---
it writes nothing --- but fails (1) because it returns a different
value every second. \code{std::printf("hi")} returns the same value
every time but fails (2) by writing to a file descriptor.

Here are some functions you have written, classified:

\begin{center}
\small
\begin{tabular}{lll}
\textbf{Function} & \textbf{Pure?} & \textbf{Why} \\
\hline
\code{int square(int n) \{ return n*n; \}}        & yes & input only, no effects \\
\code{void log(\dots)} writing to stderr           & no  & writes to a file descriptor \\
\code{int next\_id()} with a \code{static} counter & no  & reads/writes hidden state \\
\code{std::sort(begin, end)}                       & no  & mutates the range (side effect) \\
\code{Model update(Model m, Msg ev)} returning a new \code{Model} & yes & inputs in, value out \\
\end{tabular}
\end{center}

\subsection{Why anyone wants this}

Purity buys you four properties that imperative code does not have:

\begin{itemize}[leftmargin=*]
  \item \textbf{Referential transparency.} \code{update(m, Inc\{\})}
    can be replaced anywhere by its value. The compiler can cache it,
    you can paste it into a test, you can move it across threads ---
    none of those changes the meaning of the program.
  \item \textbf{Local reasoning.} To understand a pure function you
    read its body. You do not have to ask ``but what if some other
    thread mutates the field between line 3 and line 4?'' --- there
    are no fields.
  \item \textbf{Trivial testing.} The test is the function call.
    \code{ASSERT(update(\{5\}, Inc\{\}).first.count == 6)} needs no
    mock, no fixture, no setup.
  \item \textbf{Deterministic replay.} A pure function fed the same
    sequence of inputs produces the same sequence of outputs. Bug
    reports become reproducible by attaching a log of \code{Msg}s.
\end{itemize}

The price is structure: state changes are explicit, side effects are
described as values instead of done in place, and small programs
become a little more verbose. For programs that grow past a thousand
lines, the trade is plainly worth it. \maya{} bakes it in.

\begin{idea}
Pure = ``same inputs give same output, and nothing observable
happens on the side''. Pure functions can be tested, cached, moved
between threads, and replayed; impure functions can only be run.
\end{idea}

\subsection{Functional programming, briefly}

The family of languages built around pure functions is called
\textbf{functional}. Haskell, Elm, OCaml, F\#, Clojure, Scala, and
the pure subset of every modern language belong to this family. The
two big ideas are:

\begin{itemize}[leftmargin=*]
  \item \textbf{Data is immutable.} A ``modified'' value is a fresh
    value, not an in-place mutation. Old references keep seeing the
    old data.
  \item \textbf{Functions are first-class.} A function can be
    stored in a variable, passed to another function, returned from
    a function. Lambdas in \cpp{} (Chapter~\ref{ch:cpp}) are this
    feature.
\end{itemize}

\cpp{} is not a functional language --- it is a multi-paradigm one,
which is a polite way of saying ``it lets you do anything'' --- but
it has the pieces required to write functional code: \code{const}
for immutability, lambdas for first-class functions, \code{std::variant}
for sum types, \code{std::move} for cheap functional updates of
large values. \maya{} uses those pieces. It does not enforce purity
the way Haskell does; you can write impure \code{update}s if you
insist. The library is structured so that pure ones cost nothing
extra and impure ones gradually corrupt the benefits.

\begin{funfact}
Elm itself is a \emph{purely functional} language designed in 2012
for browser UIs --- Evan Czaplicki's PhD-era response to React. Its
biggest export is not the language but the pattern: the model /
update / view triangle now turns up under different names in Redux
(JavaScript), \code{bubbletea} (Go), \code{iced} (Rust),
\code{re-frame} (ClojureScript), SwiftUI's \code{Store}\dots and now
\maya{}. The idea outgrew the language.
\end{funfact}

\section{Where state usually lives}

The most common way to write a UI program is to keep state in mutable
fields scattered across whatever objects feel relevant, and to mutate
them from inside event handlers:

\begin{lstlisting}
class Counter {
    int count = 0;
public:
    void on_plus()  { ++count; redraw(); }
    void on_minus() { --count; redraw(); }
};
\end{lstlisting}

This works for tiny programs. It breaks down for the kind of program
where \maya{} is interesting --- programs with several concurrent
sources of events (keys, mouse, timers, async tasks), where you want
to undo, replay, snapshot, or test the logic in isolation. The state
is everywhere; the events come from everywhere; there is no single
place to put a breakpoint.

\section{The Elm shape}

The Elm Architecture says: pick one type to be your state (the
\textbf{Model}), pick one closed sum type to be your events (the
\textbf{Msg}), and write three functions:

\begin{itemize}[leftmargin=*]
  \item \code{init() -> Model} --- the starting state.
  \item \code{update(Model, Msg) -> Model} --- given old state and an
    event, return new state.
  \item \code{view(const Model\&) -> Element} --- given state,
    produce the UI to display.
\end{itemize}

Everything outside those three functions is the \emph{runtime}: it
reads keys, calls \code{update}, calls \code{view}, paints the result,
loops.

Here is the world's smallest example, the counter from
Chapter~\ref{ch:intro} but written in this style:

\begin{lstlisting}
struct Counter {
    struct Model { int count = 0; };
    struct Inc {}; struct Dec {}; struct Quit {};
    using Msg = std::variant<Inc, Dec, Quit>;

    static Model init() { return {}; }

    static Model update(Model m, Msg msg) {
        return std::visit(overload{
            [&](Inc)  -> Model { return {m.count + 1}; },
            [&](Dec)  -> Model { return {m.count - 1}; },
            [&](Quit) -> Model { return m; }
        }, msg);
    }

    static Element view(const Model& m) {
        return v(text(m.count) | Bold,
                 t<"[+/-] q quit"> | Dim) | pad<1> | border_<Round>;
    }
};
\end{lstlisting}

Read it twice. Notice:

\begin{itemize}[leftmargin=*]
  \item \code{Model} is a plain value. No pointers, no inheritance.
  \item \code{Msg} is a closed sum: \emph{exactly} \code{Inc},
    \code{Dec}, or \code{Quit}. If you add a fourth event tomorrow,
    every \code{std::visit} call in the program complains until you
    handle it.
  \item \code{update} is pure: same model + same message $\Rightarrow$
    same result. You can call it from a test with hand-rolled
    \code{Model}s and \code{Msg}s and check equality.
  \item \code{view} is pure too: same model $\Rightarrow$ same UI.
\end{itemize}

\begin{idea}
A pure \code{update} function is the smallest unit of testable
behaviour in a maya program. There is no mock terminal, no fake
event loop, no setup code --- just two values in, one value out.
\end{idea}

\section{But programs need to do things}

Real programs have to perform side effects. Reading a file. Setting
the terminal title. Spawning a subprocess. Quitting. If \code{update}
is pure, it can't do any of those.

The Elm answer is brilliant: have \code{update} \emph{return} the side
effects as values. They are descriptions of work, not work. The
runtime, which has permission to do impure things, interprets them.

This is the technique known as \textbf{algebraic effects} or
\emph{effects as data}. In Haskell it is the \code{IO} monad. In Rust
it is libraries like \code{tokio::Command}. In \maya{} it is
\code{Cmd<Msg>}, which is a \code{std::variant} of all the effects the
runtime knows how to interpret:

\begin{lstlisting}
template <typename Msg>
class Cmd {
public:
    struct None {};
    struct Quit {};
    struct Batch { std::vector<Cmd> cmds; };
    struct After { std::chrono::milliseconds delay; Msg msg; };
    struct SetTitle      { std::string title; };
    struct WriteClipboard { std::string content; };
    struct Task         { std::function<void(std::function<void(Msg)>)> run; };
    struct IsolatedTask { std::function<void(std::function<void(Msg)>)> run; };
    // ...
};
\end{lstlisting}

\code{update} now returns a \code{std::pair<Model, Cmd<Msg>>}:

\begin{lstlisting}
static std::pair<Model, Cmd<Msg>> update(Model m, Msg msg) {
    return std::visit(overload{
        [&](Inc)  -> std::pair<Model, Cmd<Msg>> {
            return {{m.count + 1}, Cmd<Msg>{}};
        },
        [&](Save) -> std::pair<Model, Cmd<Msg>> {
            return {m, Cmd<Msg>::write_clipboard(std::to_string(m.count))};
        },
        [&](Quit) -> std::pair<Model, Cmd<Msg>> {
            return {m, Cmd<Msg>::quit()};
        }
    }, msg);
}
\end{lstlisting}

\code{update} still does no I/O. It \emph{describes} I/O. The runtime
sees \code{WriteClipboard}, calls the OS clipboard API, and the next
call to \code{update} sees the world in a new state.

\subsection{Every \code{Cmd<Msg>} alternative, in detail}
\label{sec:elm:cmd-tour}

Let us walk the actual list. The file is
\file{maya/include/maya/core/cmd.hpp}; every alternative below is one
of the eleven cases of \code{Cmd<Msg>::Variant}.

\paragraph{\code{None}.} The empty command. Returned when
\code{update} has nothing for the runtime to do. The runtime visits
it and immediately returns.

\begin{lstlisting}
return {new_model, Cmd<Msg>::none()};
\end{lstlisting}

\paragraph{\code{Quit}.} Politely asks the runtime to exit its main
loop. The current frame finishes painting, the terminal driver tears
down (raw mode off, alt-screen out, KKP popped, cursor restored),
and \code{run<P>()} returns.

\begin{lstlisting}
return {model, Cmd<Msg>::quit()};
\end{lstlisting}

\paragraph{\code{After}.} Schedule a single \code{Msg} to be
delivered after a delay. The runtime sets a timer; when it fires,
the message is fed into \code{update} as if a user had triggered it.

\begin{lstlisting}
return {model, Cmd<Msg>::after(500ms, HideToast{})};
\end{lstlisting}

Use it for tooltips, transient notifications, debouncing. The timer
is cancelled implicitly if the runtime exits first.

\paragraph{\code{SetTitle}.} Sets the terminal window title.
Internally the runtime emits \code{ESC \textbackslash{}"]0;\dots\textbackslash{}"a}
(an OSC sequence). Outside fullscreen mode the title persists after
the program exits; inside, the alt-screen leave restores the prior
title.

\begin{lstlisting}
return {m, Cmd<Msg>::set_title("todo - " + std::to_string(n) + " items")};
\end{lstlisting}

\paragraph{\code{WriteClipboard}.} Writes a string to the OS
clipboard via OSC 52 (an escape sequence the terminal forwards to
the system clipboard). Works locally and over \code{ssh} if the
emulator opts in.

\paragraph{\code{Task}.} The escape hatch for arbitrary asynchronous
work. You hand the runtime a callable; it runs the callable on a
background worker pool, and the callable is free to dispatch any
number of \code{Msg}s back as results arrive.

\begin{lstlisting}
return {m.with(.loading = true),
        Cmd<Msg>::task([url](auto dispatch) {
            auto response = http_get(url);    // blocking, on a worker
            dispatch(Loaded{std::move(response)});
        })};
\end{lstlisting}

The \code{dispatch} parameter is itself a function: every call
funnels one \code{Msg} back through the runtime, which pushes it
through \code{update} on the main thread. Your task does not touch
the model and does not touch the screen; it only produces messages.

\paragraph{\code{IsolatedTask}.} Like \code{Task} but runs on a
dedicated detached \code{std::thread} instead of the shared pool.
Use it for work that might wedge on a blocking syscall --- a slow
filesystem mount, a network call without a timeout, a sandboxed
subprocess that might never return. A wedged \code{IsolatedTask}
leaks one thread; a wedged \code{Task} permanently consumes one of
the pool's slots and starves all subsequent tasks. The trade is
roughly 100--300\,\textmu{}s of \code{std::thread} construction per
call, paid only for genuinely hang-prone work.

\paragraph{\code{Batch}.} Combine multiple commands. The runtime
performs them all, in unspecified order.

\begin{lstlisting}
return {m, Cmd<Msg>::batch(
    Cmd<Msg>::set_title("saving\dots"),
    Cmd<Msg>::after(2s, ClearStatus{}),
    Cmd<Msg>::task([](auto dispatch) {
        write_to_disk();
        dispatch(Saved{});
    })
)};
\end{lstlisting}

Nested batches are flattened automatically and \code{None}s stripped
(see the \code{batch} static in \file{cmd.hpp}). The two convenient
forms are \code{batch(std::vector<Cmd>\{...\})} and the variadic
\code{batch(c1, c2, c3, ...)}.

\paragraph{\code{CommitScrollback} / \code{CommitScrollbackOverflow} /
\code{ForceRedraw}.} Three rendering-control commands relevant only
to inline-mode programs (Chapter~\ref{ch:inline}). Outside that
context, do not touch them. We will return to them when we cover
the inline renderer.

\subsection{Mapping \code{Cmd<Msg>} between message types}
\label{sec:elm:cmd-map}

Large programs split into components, each with their own
\code{Msg} type. A child component returns \code{Cmd<ChildMsg>}; the
parent's \code{update} returns \code{Cmd<ParentMsg>}. The
conversion is the \code{map} method:

\begin{lstlisting}
Cmd<ChildMsg>  c1 = Child::update(child_model, child_msg).second;
Cmd<ParentMsg> c2 = c1.map([](ChildMsg m) -> ParentMsg {
    return ChildEvent{std::move(m)};
});
\end{lstlisting}

\code{map} threads the conversion function through every alternative:
a \code{Cmd<ChildMsg>::Task} becomes a \code{Cmd<ParentMsg>::Task}
that wraps each \code{ChildMsg} it would dispatch with the
conversion. A \code{Quit} stays a \code{Quit}. A \code{Batch}
recurses element-wise. Read the implementation in \file{cmd.hpp};
it is one \code{std::visit} per alternative.

In category-theory terms, \code{Cmd} is a functor parameterised by
the message type, and \code{map} is the functor's lift. You do not
need to know that label; the trick is that you can refactor a
standalone component into a child of a parent program without
rewriting any of its effect-emitting code --- the parent's update
wraps the child's \code{Cmd} with \code{map}.

\section{Subscriptions: the other half}

\code{Cmd} describes effects the program \emph{starts}. There is also
the symmetric notion: events the program \emph{receives}. Keystrokes,
mouse moves, timers, async results --- these are sources of
\code{Msg}s. \maya{} models them with \code{Sub<Msg>}:

\begin{lstlisting}
template <typename Msg>
class Sub {
public:
    struct Keys     { std::function<std::optional<Msg>(Key)>   f; };
    struct Mouse    { std::function<std::optional<Msg>(Mouse)> f; };
    struct Tick     { std::chrono::milliseconds period; std::function<Msg()> f; };
    // ...
};
\end{lstlisting}

\code{subscribe(const Model\&) -> Sub<Msg>} is the fourth program
function. Like the others, it is pure: given a model, declare which
event sources are active right now. The runtime samples the
subscriptions, reads from those sources, and feeds the \code{Msg}s
back into \code{update}.

\subsection{Why subscriptions are declarative}
\label{sec:elm:why-declarative}

In an imperative GUI, you wire event sources in setup code:
``\code{button.on\_click(\dots)}, \code{timer.start(16ms)}, on
shutdown remember to \code{timer.stop()} and \code{button.disconnect()}
and don't leak the handler''. The wiring is stateful, and forgetting
to unwind it leaks event sources --- a category of bug every GUI
framework has seen.

Maya's \code{subscribe} is different: it is a pure function the
runtime calls on every frame. The returned \code{Sub} value
\emph{is} the current declaration of what event sources are active.
The runtime diffs the new \code{Sub} against the previous one and
takes care of starting and stopping the underlying listeners. You
do not call \code{start} or \code{stop}; you describe the desired
state, and the runtime makes the world match.

The model is React's hooks rules in a different language: declare
what you want, the framework reconciles. The trick is that ``what
you want'' is a value of an ordinary closed sum type, which means
the set of legal answers is bounded and visible.

\subsection{Every \code{Sub<Msg>} alternative}
\label{sec:elm:sub-tour}

From \file{maya/include/maya/app/sub.hpp}:

\paragraph{\code{None}.} Subscribe to nothing. The empty sub.

\paragraph{\code{OnKey}.} Subscribe to keyboard events. The filter
returns \code{std::optional<Msg>} --- \code{nullopt} to ignore the
event, a \code{Msg} to forward it to \code{update}. The filter is
called for every key press; only keys the application cares about
become messages.

\begin{lstlisting}
Sub<Msg>::on_key([](const KeyEvent& k) -> std::optional<Msg> {
    if (key_is(k, 'q')) return Quit{};
    if (key_is(k, '+')) return Inc{};
    return std::nullopt;          // ignore
});
\end{lstlisting}

\paragraph{\code{OnMouse}.} Same pattern, for mouse events: press,
release, drag, scroll. The filter is called for each event. See
Chapter~\ref{ch:events} for the structure of \code{MouseEvent}.

\paragraph{\code{OnResize}.} The window changed size. The function
receives the new \code{Size\{w, h\}}; it returns a \code{Msg} (not
optional --- you must produce some message, even a no-op one). The
runtime relays out the screen automatically after a resize; the
message is for application logic that depends on size (collapsing a
sidebar below 60 columns, for example).

\paragraph{\code{OnPaste}.} The user pasted multi-byte text
(bracketed paste; Chapter~\ref{ch:terminal}). The function receives
the full pasted string at once, returns a \code{Msg}. Use it where
you want paste to be one atomic event rather than a flood of
synthetic \code{Char} keys.

\paragraph{\code{Every}.} A timer. Every \code{interval} milliseconds
the runtime delivers the same \code{msg} to \code{update}. Use this
for clocks, polling, status refreshes. The interval is honoured
best-effort; missed ticks are coalesced rather than queued.

\paragraph{\code{AnimationFrame}.} A 60\,Hz tick. The runtime
delivers \code{msg} on every animation frame for as long as the sub
is returned. Drop it from the sub tree (do not include it during
the next \code{subscribe} call) to stop the animation. The cost
when no \code{AnimationFrame} sub is present is zero --- idle frames
are not rendered.

\paragraph{\code{Batch}.} Combine multiple subs. Order of activation
is unspecified; ordering of delivery follows the underlying event
sources.

\subsection{Subscribe is pure --- and that matters}

A subtle benefit of \code{subscribe} being pure: it can react to the
model.

\begin{lstlisting}
static Sub<Msg> subscribe(const Model& m) {
    auto keys = Sub<Msg>::on_key(/* always-on filter */);
    if (m.animating) {
        return Sub<Msg>::batch(keys,
            Sub<Msg>::on_animation_frame(Tick{}));
    }
    return keys;
}
\end{lstlisting}

While \code{m.animating} is true, the runtime starts a 60\,Hz tick.
When the model sets \code{animating} back to \code{false}, the next
\code{subscribe} call returns no \code{AnimationFrame}; the runtime
stops the tick. You never wrote a \code{stop()} call. The diffing
is the framework's job.

This matters for things like spinners, transient progress bars,
short-lived polling: subscribe while the work is pending, drop the
sub when it's done. Imperative code requires a flag, a cleanup path,
and a bug.

\begin{idea}
\code{subscribe(m)} returns a \emph{declaration} of which event
sources are wanted right now. The runtime diffs against the prior
declaration and starts / stops listeners. You never write
\code{stop()}.
\end{idea}

A complete Elm-style \maya{} program is therefore a struct with four
static functions:

\begin{center}
\small
\begin{tabular}{lll}
\code{init}      & $\;\to\;$ & \code{(Model, Cmd<Msg>)} \\
\code{update}    & \code{(Model, Msg)} $\;\to\;$ & \code{(Model, Cmd<Msg>)} \\
\code{view}      & \code{(const Model\&)} $\;\to\;$ & \code{Element} \\
\code{subscribe} & \code{(const Model\&)} $\;\to\;$ & \code{Sub<Msg>} \\
\end{tabular}
\end{center}

This four-tuple is the \code{Program} concept:

\begin{lstlisting}
template <typename P>
concept Program = requires(typename P::Model m, typename P::Msg msg) {
    typename P::Model;
    typename P::Msg;
    { P::init() }            -> std::convertible_to<std::pair<typename P::Model, Cmd<typename P::Msg>>>;
    { P::update(m, msg) }    -> std::convertible_to<std::pair<typename P::Model, Cmd<typename P::Msg>>>;
    { P::view(m) }           -> std::convertible_to<Element>;
    { P::subscribe(m) }      -> std::convertible_to<Sub<typename P::Msg>>;
};
\end{lstlisting}

\code{run<P>()} requires \code{Program<P>}; if your program is missing
a function or has a wrong signature, the compiler tells you which one
and why.

\section{Why this is worth the effort}

A pure-functional core has properties that an imperative one doesn't:

\begin{itemize}[leftmargin=*]
  \item \textbf{Testability without mocks.} You write
    \code{ASSERT(update(\{5\}, Inc\{\}).first == Model\{6\})}. No
    terminal, no event loop, no fixtures.
  \item \textbf{Determinism.} Given a recorded sequence of
    \code{Msg}s, you can reproduce any rendered frame. \maya{}
    leverages this for debugging.
  \item \textbf{Time travel.} Storing a list of past models is
    storing a list of plain values --- snapshot, replay, undo.
  \item \textbf{Local reasoning.} A bug in \code{update} cannot
    corrupt the renderer; a bug in the renderer cannot mutate the
    model. The interfaces between layers are values, not method
    calls.
\end{itemize}

The cost is some up-front structure. The four-function skeleton feels
heavy for a 30-line program; \maya{} provides simpler entry points
(\code{run(event\_fn, render\_fn)}) for those cases. For anything that
will live longer than an afternoon, the Elm shape is the right shape.

\section{The runtime loop, sketched}

The \code{run<P>()} runtime is one loop. Stripped of error handling
and the inline-mode branch, it looks like this:

\begin{lstlisting}
template <Program P>
void run(RunConfig cfg) {
    auto [model, cmd0] = P::init();
    Renderer r{cfg};
    interpret(cmd0);                          // perform initial effects

    while (!quit_requested()) {
        auto subs = P::subscribe(model);       // declare event sources
        Msg msg   = wait_for(subs);            // block until something arrives

        auto [next_model, next_cmd] = P::update(std::move(model), msg);
        model = std::move(next_model);
        interpret(next_cmd);                   // perform effects from update

        r.paint(P::view(model));               // re-render
    }
}
\end{lstlisting}

That is the entire shape. Everything else in \maya{}'s architecture
chapter is an embellishment: the canvas double-buffer, the SIMD diff,
the inline-mode renderer, the BG thread pool that interprets
\code{Task} commands. All in service of this loop.

\begin{theory}
\textbf{Algebraic effects} is a programming-language idea: separate
\emph{describing} what should happen from \emph{deciding how}. The
function returns a value of an effect algebra; an interpreter chooses
how to run it. The same description can be run for real, run in a
test (with a mock interpreter), or replayed from a log. Elm's
\code{Cmd} and \code{Sub} are a small, practical algebra of effects.
\end{theory}

\subsection{The loop, traced one frame at a time}
\label{sec:elm:trace}

Let's walk what happens between two consecutive frames of a running
program. Suppose the model has \code{count = 7} and the user has
just pressed \code{+}.

\paragraph{Step 0.} The runtime has already painted the previous
frame and is sitting in \code{wait\_for(subs)}. Its standard input is
open and watching for bytes.

\paragraph{Step 1.} The byte \code{0x2B} (ASCII \code{+}) arrives.
The runtime's input parser (Chapter~\ref{ch:terminal}) recognises a
plain printable key and produces \code{KeyEvent\{kind = Char, char =
'+'\}}.

\paragraph{Step 2.} The runtime walks the current \code{Sub} tree
looking for an \code{OnKey} filter. It finds one, calls it with the
key event, gets back \code{std::optional<Msg>(Inc\{\})}.

\paragraph{Step 3.} The runtime calls \code{P::update(std::move(model),
Inc\{\})}. The model is moved (not copied) into \code{update}. The
function returns \code{\{Model\{8\}, Cmd<Msg>::none()\}}.

\paragraph{Step 4.} The runtime moves the returned model into its
own \code{model} variable. It interprets the \code{Cmd}: \code{None}
means nothing to do.

\paragraph{Step 5.} The runtime calls \code{P::view(model)}. The
view function returns an \code{Element} tree (Chapter~\ref{ch:elements})
describing the new screen.

\paragraph{Step 6.} The renderer lays the tree out
(Chapter~\ref{ch:layout}), rasterises it into the back canvas
(Chapter~\ref{ch:canvas}), diffs against the front canvas
(Chapter~\ref{ch:simd}), and emits exactly the bytes for the cells
that changed (Chapter~\ref{ch:ansi}). The front and back canvases
swap.

\paragraph{Step 7.} The runtime calls \code{P::subscribe(model)} to
learn what event sources to watch this frame, diffs against the
prior subscription, and re-enters \code{wait\_for(subs)}.

The whole cycle is the loop body. Every keystroke, every timer tick,
every async result re-runs steps 1--7. There is no other entry
point into your program's logic.

\section{Testing the model in isolation}
\label{sec:elm:tests}

The reward for keeping \code{update} pure is that you can test the
business logic without spinning up a terminal at all. There is no
mock event loop, no fake renderer, no setup of fixtures. The test is
a function call.

\begin{lstlisting}
#include <cassert>
#include "counter.hpp"

int main() {
    using M = Counter::Model;
    auto [m0, _0] = Counter::init();
    assert(m0.count == 0);

    auto [m1, _1] = Counter::update(m0, Counter::Inc{});
    auto [m2, _2] = Counter::update(m1, Counter::Inc{});
    auto [m3, _3] = Counter::update(m2, Counter::Inc{});
    assert(m3.count == 3);

    auto [m4, _4] = Counter::update(m3, Counter::Dec{});
    assert(m4.count == 2);

    auto [m5, c5] = Counter::update(m4, Counter::Quit{});
    assert(c5.is_quit());
}
\end{lstlisting}

No terminal is opened. No bytes are written. The program prints
nothing and either succeeds or aborts on a failed assertion. Compile
it with whatever \cpp{} testing framework you like ---
\code{Catch2}, \code{doctest}, GoogleTest, or none at all
(\code{assert} is fine for small programs).

\subsection{Testing effects}

What if \code{update} returns a non-trivial \code{Cmd}? You inspect
the returned variant:

\begin{lstlisting}
auto [m, cmd] = MyApp::update(m_in, Save{});
assert(std::holds_alternative<Cmd<Msg>::WriteClipboard>(cmd.inner));
auto& w = std::get<Cmd<Msg>::WriteClipboard>(cmd.inner);
assert(w.content == "expected text");
\end{lstlisting}

The \code{Cmd} is a value; you do not have to perform it to check it.
This is the deepest practical benefit of effects-as-data: the test
asks ``what would the program ask the runtime to do?'' without
asking the runtime to do anything.

For \code{Task} commands the situation is trickier --- the inside of
the task is a \code{std::function} that you would normally have to
run to find out what messages it produces. Two practical paths:

\begin{enumerate}[leftmargin=*]
  \item Make the task call a small helper function (e.g.
    \code{fetch\_user(int)}) instead of inline code, and test the
    helper separately.
  \item Run the task synchronously in a test: hand it a lambda that
    records every \code{Msg} into a vector, then assert on the
    vector.
\end{enumerate}

\begin{lstlisting}
std::vector<Msg> dispatched;
auto run_task = [&](auto& task_cmd) {
    task_cmd.run([&](Msg m) { dispatched.push_back(std::move(m)); });
};
\end{lstlisting}

In maya's own test suite, both patterns appear. The library's
internal HTTP layer factors out \code{http\_get\_blocking} so the
\code{Cmd::task} wrapping it can be tested independently.

\subsection{Property tests}

Because \code{update} is pure, property-based testing is a natural
fit. \emph{For every model \code{m} and every \code{Inc} message,
the resulting count is one more than \code{m.count}.} That is one
\code{rapidcheck} property:

\begin{lstlisting}
rc::check("inc increments by one", [](Counter::Model m) {
    auto [m1, _] = Counter::update(m, Counter::Inc{});
    RC_ASSERT(m1.count == m.count + 1);
});
\end{lstlisting}

The property generator produces hundreds of random models and
verifies the invariant on each. Imperative GUI code typically cannot
be property-tested at all --- the state is too entangled.

\section{Replay, snapshot, time travel}
\label{sec:elm:replay}

If you record every \code{Msg} that arrives, you can reproduce any
state the program ever reached by replaying the log into
\code{update}, starting from \code{init}. This is not a research
property; it is the natural consequence of \code{update} being pure.

A skeletal recorder, wrapping any \code{Program}:

\begin{lstlisting}
template <Program P>
struct Recorded {
    using Model = std::pair<typename P::Model, std::vector<typename P::Msg>>;
    using Msg   = typename P::Msg;

    static auto init() -> std::pair<Model, Cmd<Msg>> {
        auto [m, c] = P::init();
        return {{std::move(m), {}}, c};
    }
    static auto update(Model state, Msg msg) -> std::pair<Model, Cmd<Msg>> {
        auto& [inner, log] = state;
        log.push_back(msg);                  // record
        auto [m1, c] = P::update(std::move(inner), msg);
        return {{std::move(m1), std::move(log)}, c};
    }
    static Element view(const Model& state) { return P::view(state.first); }
    static Sub<Msg> subscribe(const Model& state) {
        return P::subscribe(state.first);
    }
};
\end{lstlisting}

The recorded log is just \code{std::vector<Msg>}. To replay, fold
\code{update} over the log starting from a fresh \code{init}:

\begin{lstlisting}
auto [m, _] = P::init();
for (auto& msg : recorded_log) {
    auto [next, _] = P::update(std::move(m), msg);
    m = std::move(next);
}
// m is now exactly what the user saw at the end of the recorded session
\end{lstlisting}

This is the heart of every ``undo'' implementation: store enough
messages to rewind, replay the prefix. \emph{Redux DevTools} does the
same for JavaScript apps; Elm's debugger likewise. Once your update
is pure, you get this for free.

\begin{idea}
Log the \code{Msg}s. Replay through \code{update}. Get the same
model back, every time. This is undo, snapshot, debug-record, and
deterministic bug repro --- one fact about pure functions, four
product features.
\end{idea}

\section{Splitting a big program into components}
\label{sec:elm:components}

A single \code{Model} / \code{Msg} pair works fine up to a few
thousand lines. Beyond that, you want to split: a child component
has its own \code{Model} and \code{Msg}, and the parent embeds them.

The machinery is mechanical:

\begin{lstlisting}
struct Search {
    struct Model { std::string query; std::vector<std::string> results; };
    struct Type { char c; };
    struct Backspace {};
    struct ResultsArrived { std::vector<std::string> r; };
    using Msg = std::variant<Type, Backspace, ResultsArrived>;

    static auto update(Model, Msg) -> std::pair<Model, Cmd<Msg>>;
    static Element view(const Model&);
    static Sub<Msg> subscribe(const Model&);
};

struct App {
    struct Model {
        Search::Model search;
        // ... other components ...
    };
    struct SearchEvent { Search::Msg child; };
    struct GlobalQuit {};
    using Msg = std::variant<SearchEvent, GlobalQuit>;

    static auto update(Model m, Msg msg) -> std::pair<Model, Cmd<Msg>> {
        return std::visit(overload{
            [&](GlobalQuit) -> std::pair<Model, Cmd<Msg>> {
                return {m, Cmd<Msg>::quit()};
            },
            [&](SearchEvent ev) -> std::pair<Model, Cmd<Msg>> {
                auto [child_model, child_cmd] =
                    Search::update(std::move(m.search), ev.child);
                m.search = std::move(child_model);
                auto parent_cmd = child_cmd.map([](Search::Msg cm) {
                    return Msg{SearchEvent{std::move(cm)}};
                });
                return {std::move(m), std::move(parent_cmd)};
            },
        }, msg);
    }
};
\end{lstlisting}

What to notice:

\begin{itemize}[leftmargin=*]
  \item The child's \code{Model} is a field of the parent's
    \code{Model}. No pointers, no inheritance.
  \item The child's \code{Msg} is wrapped by a parent message
    constructor (\code{SearchEvent}) into the parent's \code{Msg}.
  \item When the parent's \code{update} sees a \code{SearchEvent},
    it forwards to the child's \code{update}, getting back
    \code{(ChildModel, Cmd<ChildMsg>)}.
  \item The new child model is written back into the parent model.
    The child \code{Cmd} is mapped into a parent \code{Cmd} with
    \code{.map(\dots)}.
\end{itemize}

\code{subscribe} composes the same way: the parent's
\code{subscribe} calls the child's, wraps the result with
\code{Sub<>::map(\dots)}, and returns it. \code{view} embeds the
child's view as a subtree.

This is the entire technique for scaling Elm-style code. There is no
``component framework''; it is plain \cpp{} types composed by
hand. The compiler keeps everyone honest about which messages can
appear where.

\subsection{When to split, when to keep it flat}

Resist the urge to split too early. A flat 2,000-line program with
clear \code{Msg} categories is easier to read than three components
wrapped in two layers of \code{map} indirection. Split when:

\begin{itemize}[leftmargin=*]
  \item A subtree of the UI is genuinely reusable (a date picker,
    an autocomplete) and you want to drop it into multiple
    programs.
  \item The \code{Msg} variant has more than 20 alternatives and the
    \code{update} function visibly groups into clusters.
  \item A subteam owns a feature and you want their changes to land
    in one file, not in a global \code{Msg}.
\end{itemize}

Otherwise, flat is fine. Maya's own programs in \code{agentty} are
mostly flat with a few well-isolated child components for
specifically reusable pieces.

\section{A worked refactor: from script to component}
\label{sec:elm:worked-refactor}

For concreteness, here is a small but realistic program in both the
quick-tool form and the Program form. The program is a to-do list:
the user types items into a draft buffer, presses \code{Enter} to
commit, presses \code{space} to toggle the selected item's done
state, \code{d} to delete, arrow keys to navigate, \code{q} to quit.

First the imperative sketch, the kind of code you might write in an
afternoon and regret later:

\begin{lstlisting}
struct Item { std::string text; bool done = false; };

int main() {
    std::vector<Item> items;
    std::string draft;
    int selected = 0;

    run({.title = "todo"},
        [&](const Event& ev) {
            if (key(ev, 'q')) return false;
            if (key(ev, '\n')) {
                if (!draft.empty()) {
                    items.push_back({draft, false});
                    draft.clear();
                }
            } else if (key(ev, ' ')) {
                if (!items.empty()) items[selected].done = !items[selected].done;
            } else if (key(ev, 'd')) {
                if (!items.empty()) {
                    items.erase(items.begin() + selected);
                    if (selected >= (int)items.size())
                        selected = std::max(0, (int)items.size() - 1);
                }
            } else if (ev.is_key(KeyKind::Up))   selected = std::max(0, selected - 1);
            else if   (ev.is_key(KeyKind::Down)) selected = std::min((int)items.size() - 1,
                                                                    selected + 1);
            else if   (ev.is_key(KeyKind::Backspace)) {
                if (!draft.empty()) draft.pop_back();
            } else if (ev.is_char()) draft.push_back(ev.character());
            return true;
        },
        [&] {
            // ... build the view from items, draft, selected ...
        });
}
\end{lstlisting}

Works. Two issues: state lives in main's stack frame (no test can
reach it without launching the program), and the event-handler
branches are not enumerable --- adding a feature means adding
another \code{if} and hoping you covered every case.

Now the same program as a \code{Program}:

\paragraph{Step 1: define the \code{Model}.}

\begin{lstlisting}
struct Todos {
    struct Item { std::string text; bool done = false; };
    struct Model {
        std::vector<Item> items;
        std::string       draft;
        int               selected = 0;
    };
};
\end{lstlisting}

A plain value type. Equality is comparison of fields. Copying is a
vector copy plus two strings.

\paragraph{Step 2: enumerate the \code{Msg}.}

\begin{lstlisting}
    struct AppendDraft   { char c; };
    struct DeleteDraft   {};
    struct SubmitDraft   {};
    struct ToggleDone    {};
    struct DeleteCurrent {};
    struct SelectUp      {};
    struct SelectDown    {};
    struct Quit          {};
    using Msg = std::variant<AppendDraft, DeleteDraft, SubmitDraft,
                              ToggleDone, DeleteCurrent, SelectUp,
                              SelectDown, Quit>;
\end{lstlisting}

Eight cases. The compiler will refuse to compile an \code{update}
that does not handle all of them.

\paragraph{Step 3: write \code{update}.}

\begin{lstlisting}
    static auto update(Model m, Msg msg) -> std::pair<Model, Cmd<Msg>> {
        return std::visit(overload{
            [&](AppendDraft a) -> std::pair<Model, Cmd<Msg>> {
                m.draft.push_back(a.c);
                return {std::move(m), Cmd<Msg>::none()};
            },
            [&](DeleteDraft) -> std::pair<Model, Cmd<Msg>> {
                if (!m.draft.empty()) m.draft.pop_back();
                return {std::move(m), Cmd<Msg>::none()};
            },
            [&](SubmitDraft) -> std::pair<Model, Cmd<Msg>> {
                if (!m.draft.empty()) {
                    m.items.push_back({std::move(m.draft), false});
                    m.draft.clear();
                }
                return {std::move(m), Cmd<Msg>::none()};
            },
            [&](ToggleDone) -> std::pair<Model, Cmd<Msg>> {
                if (!m.items.empty())
                    m.items[m.selected].done = !m.items[m.selected].done;
                return {std::move(m), Cmd<Msg>::none()};
            },
            [&](DeleteCurrent) -> std::pair<Model, Cmd<Msg>> {
                if (!m.items.empty()) {
                    m.items.erase(m.items.begin() + m.selected);
                    if (m.selected >= (int)m.items.size())
                        m.selected = std::max(0, (int)m.items.size() - 1);
                }
                return {std::move(m), Cmd<Msg>::none()};
            },
            [&](SelectUp) -> std::pair<Model, Cmd<Msg>> {
                m.selected = std::max(0, m.selected - 1);
                return {std::move(m), Cmd<Msg>::none()};
            },
            [&](SelectDown) -> std::pair<Model, Cmd<Msg>> {
                m.selected = std::min((int)m.items.size() - 1, m.selected + 1);
                return {std::move(m), Cmd<Msg>::none()};
            },
            [&](Quit) -> std::pair<Model, Cmd<Msg>> {
                return {std::move(m), Cmd<Msg>::quit()};
            },
        }, msg);
    }
\end{lstlisting}

Every case is named. Adding a ninth feature is adding a struct to
the variant and a lambda to the visit; the compiler will list every
callsite that needs updating.

\paragraph{Step 4: write \code{view}.}

\begin{lstlisting}
    static Element view(const Model& m) {
        std::vector<Element> rows;
        rows.reserve(m.items.size());
        for (int i = 0; i < (int)m.items.size(); ++i) {
            const auto& it = m.items[i];
            auto check  = it.done ? "[x] " : "[ ] ";
            auto style  = it.done ? (Style{} | Dim | Strike) : Style{};
            auto line   = text(std::string(check) + it.text) | style;
            if (i == m.selected) line = line | Inverse;
            rows.push_back(line);
        }
        auto list  = rows.empty() ? Element{t<"(empty)"> | Dim} : v(rows);
        auto input = h(t<"> "> | Bold,
                       text(m.draft + "_") | Fg<200, 200, 100>);
        auto help  = t<"enter add | space toggle | d delete | q quit"> | Dim;

        return v(list | pad<1> | border_<Round>,
                 input | pad<1>,
                 help) | pad<1>;
    }
\end{lstlisting}

The view depends only on the model. Hand it any \code{Model} value
and inspect the returned \code{Element} tree without rendering it.

\paragraph{Step 5: write \code{subscribe}.}

\begin{lstlisting}
    static Sub<Msg> subscribe(const Model&) {
        return Sub<Msg>::on_key([](const KeyEvent& k) -> std::optional<Msg> {
            if (key_is(k, 'q'))                 return Quit{};
            if (k.kind == KeyKind::Enter)       return SubmitDraft{};
            if (k.kind == KeyKind::Backspace)   return DeleteDraft{};
            if (k.kind == KeyKind::Up)          return SelectUp{};
            if (k.kind == KeyKind::Down)        return SelectDown{};
            if (key_is(k, ' '))                 return ToggleDone{};
            if (key_is(k, 'd'))                 return DeleteCurrent{};
            if (k.kind == KeyKind::Char)        return AppendDraft{k.character};
            return std::nullopt;
        });
    }
\end{lstlisting}

The keymap is one screen of code. Every other piece of input
handling in the program goes through this single filter.

\paragraph{Step 6: \code{init} and \code{main}.}

\begin{lstlisting}
    static auto init() -> std::pair<Model, Cmd<Msg>> {
        return {Model{}, Cmd<Msg>::none()};
    }
}; // end struct Todos

int main() { run<Todos>({.title = "todo"}); }
\end{lstlisting}

The whole application is one file, a few hundred lines, and every
part is nameable. Tests look like:

\begin{lstlisting}
int main() {
    using T = Todos;
    auto [m0, _] = T::init();

    auto [m1, _1] = T::update(std::move(m0), T::AppendDraft{'h'});
    auto [m2, _2] = T::update(std::move(m1), T::AppendDraft{'i'});
    auto [m3, _3] = T::update(std::move(m2), T::SubmitDraft{});
    assert(m3.items.size() == 1 && m3.items[0].text == "hi");
    assert(m3.draft.empty());

    auto [m4, _4] = T::update(std::move(m3), T::ToggleDone{});
    assert(m4.items[0].done);

    auto [m5, _5] = T::update(std::move(m4), T::DeleteCurrent{});
    assert(m5.items.empty());
}
\end{lstlisting}

\section{Performance: where the loop spends time}
\label{sec:elm:perf}

The loop is simple, but ``one frame per event'' could be expensive
if executed naively. A few notes on what \maya{} does to keep it
fast and where the costs actually fall.

\paragraph{The model is moved, not copied.} \code{update(Model,
Msg)} takes the model \emph{by value}; the runtime
\code{std::move}s the current model in. For a \code{Model} of small
fields, this is the cost of a few pointer copies; for one with a
big vector or string, the buffer's ownership swaps and no
element-by-element copy happens.

\paragraph{The view is rebuilt every frame.} The \code{Element}
tree is allocated fresh on each frame. This sounds wasteful but is
cheap: small bump-style allocations, freed when the tree goes out of
scope at the end of the frame. Chapter~\ref{ch:canvas} shows the
rendering pipeline that absorbs this cost into a single linear
sweep.

\paragraph{Diff, not full repaint.} The renderer compares the new
canvas to the previous one and emits bytes only for changed cells.
A frame with no visible changes emits zero bytes. A frame whose only
change is one cell emits five to ten bytes. The SIMD pass in
Chapter~\ref{ch:simd} compares 32 cells per CPU instruction; a 200x60
terminal compares in microseconds.

\paragraph{Idle frames have no cost.} If no event arrives and no
subscription is firing, the runtime is blocked in a \code{poll}
syscall and uses zero CPU. There is no render-on-tick busy loop. A
\maya{} program that nobody is using is invisible in \code{htop}.

\paragraph{Tasks run on a worker pool.} \code{Cmd::task} dispatches
to a small bounded pool of background threads; \code{Cmd::task\_isolated}
spins a fresh thread per call. The main loop never blocks on user
code. A wedged task does not freeze the UI.

\begin{recap}
The loop is cheap because: the model is moved not copied; the
element tree is allocated cheaply and dropped per-frame; the diff
is SIMD-fast; idle frames cost zero; tasks are off the main thread.
A \maya{} program redraws at hundreds of frames per second on a
modern laptop, and at dozens over slow ssh links.
\end{recap}

\section{Pitfalls and anti-patterns}
\label{sec:elm:pitfalls}

If you ignore Elm's discipline, the library still compiles and runs,
but you start losing the benefits. The common ways:

\begin{warning}
\textbf{Hidden mutation in update.} Calling
\code{some\_global.cache.set(\dots)} inside \code{update} works but
makes the function impure; testing becomes ``set up the global,
then call, then check the global''. Push the mutation out as a
\code{Cmd::task} that performs it, or move the state into the
model.
\end{warning}

\begin{warning}
\textbf{Synchronous I/O in update.} \code{std::fopen} inside
\code{update} blocks the main loop. The UI freezes. Always wrap
file/network/subprocess work in a \code{Cmd::task} so it runs on a
worker.
\end{warning}

\begin{warning}
\textbf{Coupling view to side effects.} \code{view} should depend
only on the model. If it reads the clock, the filesystem, or any
mutable global, frames stop being deterministic, and the diff can
produce a flicker between two equally valid renderings of ``now''.
Capture the world's current value into the model via a
subscription (\code{Sub::every}) and read it from the model.
\end{warning}

\begin{warning}
\textbf{Overusing \code{Cmd::task} for trivial work.} A task is
cheap but not free --- pool dispatch, worker wake, mutex hop. If
your ``async'' work is \code{counter + 1}, just compute it inline.
\end{warning}

\begin{warning}
\textbf{Forgetting the \code{std::optional} in \code{OnKey}.} The
filter must return \code{std::nullopt} for keys the application
doesn't care about. Returning a \code{Msg} for every key produces a
flood and may interfere with the runtime's own bindings (\code{Ctrl-C}
in some configurations).
\end{warning}

\section{When to skip the Program scaffold}
\label{sec:elm:simple}

Elm-style discipline pays for itself on programs that grow.
For a one-shot script --- print a banner, ask one question, exit ---
it is overkill. \maya{} provides smaller entry points for those
cases.

\code{maya::print(element)} renders a static \code{Element} once and
returns. No model, no loop, no subscriptions. We used it for the
five-line banner in Chapter~\ref{ch:intro}.

The second-smallest form, useful when you need a single keystroke
but nothing more elaborate:

\begin{lstlisting}
std::println("press y to continue");
auto key = maya::read_key();
if (key == 'y') /* go */;
\end{lstlisting}

For anything with state that changes over time, switch to the
\code{Program} form. The boundary is roughly: if your program has
three distinct keys it responds to, write the \code{Program}.

\section{From single-function to Program: a worked refactor}

Here is the same program in both styles. First, the quick-tool form
(useful for one-off scripts):

\begin{lstlisting}
Signal<int> n{0};
run({.title = "counter"},
    [&](const Event& ev) {
        if (key(ev, '+')) n.update([](int& x){ ++x; });
        if (key(ev, '-')) n.update([](int& x){ --x; });
        return !key(ev, 'q');
    },
    [&] {
        return v(text("Count: " + std::to_string(n.get())) | Bold,
                 t<"[+/-] q quit"> | Dim) | pad<1>;
    });
\end{lstlisting}

This is fine. But the state is captured by reference inside a lambda,
the event handler returns a \code{bool}, and there is no separation
between business logic and presentation.

The Program version:

\begin{lstlisting}
struct Counter {
    struct Model { int count = 0; };
    struct Inc {}; struct Dec {}; struct Quit {};
    using Msg = std::variant<Inc, Dec, Quit>;

    static auto init() -> std::pair<Model, Cmd<Msg>> {
        return {{}, Cmd<Msg>::none()};
    }
    static auto update(Model m, Msg msg) -> std::pair<Model, Cmd<Msg>> {
        return std::visit(overload{
            [&](Inc)  { return std::pair{Model{m.count + 1}, Cmd<Msg>{}}; },
            [&](Dec)  { return std::pair{Model{m.count - 1}, Cmd<Msg>{}}; },
            [](Quit)  { return std::pair{Model{}, Cmd<Msg>::quit()}; }
        }, msg);
    }
    static Element view(const Model& m) {
        return v(text("Count: " + std::to_string(m.count)) | Bold,
                 t<"[+/-] q quit"> | Dim) | pad<1>;
    }
    static Sub<Msg> subscribe(const Model&) {
        return key_map<Msg>({{'+', Inc{}}, {'-', Dec{}}, {'q', Quit{}}});
    }
};
int main() { run<Counter>({.title = "counter"}); }
\end{lstlisting}

It is a few more lines but the parts are now nameable. Tests look
like:

\begin{lstlisting}
TEST_CASE("incrementing the counter") {
    auto [m, _] = Counter::update({.count = 5}, Counter::Inc{});
    CHECK(m.count == 6);
}
\end{lstlisting}

No terminal, no event loop, no render. That is the entire point.

\begin{recap}
A \maya{} application is a \code{Program}: a \code{Model} type, a
\code{Msg} sum, and four pure functions \code{init}/\code{update}/
\code{view}/\code{subscribe}. Effects (\code{Cmd}) and event sources
(\code{Sub}) are described as data; the runtime interprets them. The
result is a program whose business logic is trivially testable.
\end{recap}

\section*{Workshop}

\begin{enumerate}[leftmargin=*]
  \item Add a \code{Reset} message to the counter that sets
    \code{count} back to zero. Add a \code{r} keybinding.
  \item Add a \code{Tick} subscription that increments the counter
    once per second.
  \item Add a \code{Save} message that writes the current count to
    the clipboard via \code{Cmd<Msg>::write\_clipboard(...)}.
  \item Write a unit test (just a \code{static\_assert} or a
    one-line \code{main}) that pumps three \code{Inc}s through
    \code{update} and checks the final count.
\end{enumerate}
